\chapter*{译校说明}
\addcontentsline{toc}{chapter}{译校说明}
\markboth{译校说明}{}

\section*{一、底本与参照本}

本书底本是 János Kollár 与森重文合著的日文专著『双有理幾何学』（岩波講座 現代数学の展開，岩波書店，1998）。该日文本并非独立著作，而是两位作者 1998 年由剑桥大学出版社出版的英文专著 \emph{Birational Geometry of Algebraic Varieties}（Cambridge Tracts in Mathematics 134）的日译本，章节结构与英文版一一对应（绪论、七章、参考文献、索引）。日译本在术语与行文上带有日本数学界的习惯，本译注本以日文本为唯一底本，逐页照译；英文原著仅作语义参照，用于校核长句结构、专名与文献，凡两者有实质出入处，一律以日文本为准，并在译注中标明。

需要提醒读者：本书初版于 1998 年，正文中的活动区域（如极小模型纲领的四维以上情形、翻转的存在性、广义对数终端奇点的收缩定理等）在 1998 年之后已有重大进展。译文不作改写，只在章末评注与“研读札记”中指出这些地方的后续发展与今天的读法。

日文本卷首的作者前言（はしがき，原书 p.v–vi）与谢辞一并译出，置于卷首；前言中提到的“英文版缺少命题 2.71 与命题 4.11(2)”一语，可在两版对照阅读时作索引。

\section*{二、体例}

\begin{enumerate}[leftmargin=2.2em,itemsep=2pt]
\item \textbf{页码对应}：每页正文前以 \verb|% [原文 p.N]| 标出日文底本页码；页边不印页码对照，正文中的“原书 p.N”均指日文本页码。
\item \textbf{编号}：章、节、定理、命题、引理、推论、定义、例、注、公式的编号一律照抄日文本，不重编、不跳号。正文交叉引用一律用原编号的纯文本形式，如“定理 1.4”“式 (1.12)”。
\item \textbf{陈述环境的尾符}：日文本在“绪论”第 (16) 条中规定：证明末尾记■；定义、例等的末尾记□；定理、命题等，不给出证明时记■，证明在后面给出时记□。本译注本沿用此规则，因此读者会看到与常见中文教材相反的用法——定理末尾的■表示“本书不证”，□表示“证明在下文”。
\item \textbf{尾符的第三档}：实际排印时还有第三档——证明紧跟在声明之后时，声明末尾什么符号都不印，■ 只出现在证明末尾。书中这类声明并不少见（例如引理 3.15 之后的证明紧接其后）。本译注本对这一档一律照原书排印：声明末尾留空、不作任何标记，请读者不要据此误判为“未证”。
\item \textbf{批注}：译注分三类：术语与记号提醒、原书省略步骤的补足、背景与直觉。批注不改变正文，在 LaTeX 源码中以 \verb|% [Tutor Note: …]| 形式嵌在相应位置，不排入正文；全部译注汇编于书末附录 B（按章、按原书页码排列），正文中相应位置不另作标记。
\item \textbf{本节重点与章末评注}：每节开头的“本节重点”概括该节的对象、工具与结论；每章末尾的“本章回顾与评注”梳理全章脉络、核心结论与易误之处。二者均为译者所加，非原书内容。
\item \textbf{图}：原书插图按图题编号；纯线条示意图按原图重绘、图内文字改用中文，其余照原样裁自页图并在图题中注明。
\item \textbf{索引}：原书末尾的索引（原书 p.311 起）按“术语（日、英）—中译”的体例不予逐条翻译，检索功能改由两个附录承担——附录 A“术语对照表”按英文、日文查中译，附录 D“定理与定义索引”按编号查类型与原书页码。
\item \textbf{附录}：附录 A 为术语对照表；附录 B 汇编正文的全部译注；附录 C 为译者本人的“研读札记”（按章记录理解路径、技术要点与未解疑问，与各章末“本章回顾与评注”分工不同）；附录 D 为定理与定义索引（原书无此附录，为检索便利所加）。
\end{enumerate}

\section*{三、术语}

术语以书末附录“术语对照表（日—中—英）”为准。若干译名需先说明：

\begin{description}[leftmargin=2.2em,style=nextline,itemsep=1pt]
\item[terminal / canonical / log terminal / log canonical] 依次译作\textbf{终端}、\textbf{标准}、\textbf{对数终端}、\textbf{对数标准}。日文本作“端末”“標準”“ログ端末”“ログ標準”。“终端奇点”与“标准奇点”是极小模型纲领的两类基本奇点，与“极射线（extremal ray）”不是一回事，请读者注意区分。
\item[extremal ray] 译作\textbf{极射线}。日文本作“端射線”，字面易与“端末（terminal）”混淆，故中文取“极”字。
\item[flip / flop] 依次译作\textbf{翻转}、\textbf{翻模}。二者是三维极小模型纲领中处理小收缩的两类操作，互为对偶，英文原词形近而义别，中文亦取形近的两个词以便对照。
\item[Q-factorial / Q-factorialization] 译作\textbf{Q-分解的}、\textbf{Q-分解化}。日文本作“Q-分解”，系英文术语的日译缩略，中文照日文形制，首次出现处加译注说明其含义（每个 Weil 除子都是 Q-Cartier 除子）。
\item[inversion of adjunction] 译作\textbf{伴随反转}。日文本作“逆同伴”，字面为“逆伴随”，中文取“伴随反转”以明其义。
\item[vanishing theorem] 译作\textbf{消没定理}（不用“消去”“消灭”），“non-vanishing”译“非消没”。
\item[deformation] 译作\textbf{形变}（不用“变形”）。
\end{description}

\section*{四、记号}

除原书记号外，本译注本统一：$\mathbb{P}^n$ 为 $n$ 维射影空间，$\mathcal{O}_X$ 为结构层，$K_X$ 为标准除子（标准类），$\rho(X)$ 为 Picard 数，$\overline{NE}(X)$ 为曲线锥的闭包，$\kappa(X)$ 为小平维数。层用花体（$\mathcal{O},\mathcal{I},\mathcal{L}$），理想层用 $\mathcal{I}$；除子用斜体大写（$D,H,K$）。

\section*{五、致读者}

这是一本“从曲线与曲面的经典理论通向高维极小模型纲领”的书，密度大、跳跃多，作者的写法是“给出关键想法、把技术细节留给文献”。读第一遍时不妨先读每节的“本节重点”与每章的“本章回顾与评注”，再回头逐句读正文；书末的“研读札记”记录了译者在逐章翻译过程中的理解路径与卡壳处，或许对同样在啃这本书的读者有些用处。

译文中的错误由译者负责，欢迎指正。

\vspace{1em}
\noindent\hfill 译者\quad 二〇二六年十一月
